\documentclass[11pt,letterpaper]{article}

\usepackage[margin=1in]{geometry}
\usepackage{setspace}
\usepackage{booktabs}
\usepackage{hyperref}
\usepackage{enumitem}
\usepackage{titlesec}
\usepackage{xcolor}

\hypersetup{
  colorlinks=true,
  linkcolor=black,
  urlcolor=blue,
  citecolor=black
}

\setstretch{1.12}
\setlist[itemize]{leftmargin=1.4em,itemsep=0.25em,topsep=0.3em}
\setlist[enumerate]{leftmargin=1.4em,itemsep=0.25em,topsep=0.3em}
\titleformat{\section}{\large\bfseries}{\thesection}{0.6em}{}
\titleformat{\subsection}{\normalsize\bfseries}{\thesubsection}{0.6em}{}

\title{\textbf{CAI6734 Week~6 Individual Progress Report}\\[0.35em]
\large Testing in Imagination:\\
World-Model Planning for Language Agents\\
in Sequential Clinical Diagnosis}
\author{Daniel Escalante\\[0.4em]
\small Course: CAI6734 (University of Florida)\\
Advisor: Xuefeng Liu}
\date{October 6, 2026}

\begin{document}
\maketitle
\thispagestyle{empty}

\noindent
This individual progress report covers my work for Week~6 on the Testing in Imagination / LA-CDM reproduction project. I write in the first person and report only verified progress---no fabricated metrics or PHI.

\section{Progress this week}

\subsection{CITI / PhysioNet prerequisites}
On 2026-10-06 I completed the CITI \emph{Human Research -- Data or Specimens Only Research} course (MIT Affiliates; Record ID 80147027; score 100\%). This is a PhysioNet prerequisite. I have \textbf{not} yet downloaded any PhysioNet dataset.

I also started PhysioNet credentialing. My supervisor reference is Xuefeng Liu (\texttt{xuefeng.liu@ufl.edu}). The research-topic text I am using describes reproduction of MIMIC-IV-Ext CDM / LA-CDM for CAI6734 (Testing in Imagination).

\subsection{LA-CDM reproduction scaffold}
I built the \texttt{lacdm/} reproduction scaffold in our course repository:
\begin{itemize}
  \item Upstream LA-CDM pinned as a submodule at commit \texttt{3f435a10}.
  \item Nemotron Hydra overrides for the CoC-compliant backbone (logged deviation D-5): \texttt{nvidia/Llama-3.1-Nemotron-Nano-8B-v1}.
  \item Phase~E unit tests: \textbf{83 passed} on my Mac (no CUDA required).
\end{itemize}

\subsection{Git history cleanup and merge}
I resolved a temporary local orphan / unrelated-histories situation by porting the \texttt{lacdm} content onto \texttt{origin/main} rather than merging unrelated histories. That work landed on \texttt{main} at commit \texttt{95ce5fc}.

\paragraph{Local working path.}
\texttt{Documents/AIBHS/Applied Generative AI in Medicine/Testing in Imagination}
(private repo: \url{https://github.com/scalinity/testing-in-imagination}).

\section{Setbacks and how I addressed them}

\begin{itemize}
  \item \textbf{Full Table~1 / GRPO blocked.} My Mac has no CUDA; PhysioNet credentialing is still pending; I need a HiPerGator path before any real GRPO or vLLM job. I kept progress on Mac-OK work (scaffold, Hydra overrides, Phase~E tests) and documented the GPU/data gate rather than inventing run numbers.
  \item \textbf{Temporary local orphan git history vs GitHub.} Early local history diverged from \texttt{origin/main}. I fixed this by porting content onto \texttt{origin/main} (feat branch then merge) instead of an unrelated-histories merge, so \texttt{main} stays a single linear project history ending at \texttt{95ce5fc}.
\end{itemize}

\section{Next steps}

\begin{enumerate}
  \item Finish PhysioNet credentialing and request MIMIC-IV-Ext CDM; confirm authorization with Edwin and Hugo before any data prep that could involve PHI.
  \item Pin Nemotron on HiPerGator; override any remaining Qwen Hydra defaults before the first cluster job (D-5).
  \item After data access: run S2 \texttt{prepare\_data} / smoke eval; meanwhile continue the synthetic CDM path under \texttt{extension/}.
\end{enumerate}

\section*{Notes}
No Table~1 accuracy, F1, or cost figures are reported this week---those require PhysioNet-authorized data and CUDA training/eval that I have not run. Repository stewardship and Mac-side Phase~E verification are the deliverables I completed.

\end{document}
